\begin{frame}{Definição}
\begin{block}{Registrador de deslocamento}
É um circuito sequencial formado por flip-flops, capaz de armazenar bits e
mover esses bits de uma posição para outra a cada pulso de clock.
\end{block}

\begin{center}
\begin{tikzpicture}[>=Latex, node distance=0.35cm, ff/.style={draw, minimum width=0.9cm, minimum height=1cm}]
  \node (in) {Entrada serial};
  \node[ff, right=of in] (q0) {$Q_0$};
  \node[ff, right=of q0] (q1) {$Q_1$};
  \node[ff, right=of q1] (q2) {$Q_2$};
  \node[ff, right=of q2] (q3) {$Q_3$};
  \node[right=of q3] (out) {Saída serial};
  \node[below=0.8cm of q1] (clk) {Clock comum};
  \draw[->, thick] (in) -- (q0);
  \draw[->, thick] (q0) -- (q1);
  \draw[->, thick] (q1) -- (q2);
  \draw[->, thick] (q2) -- (q3);
  \draw[->, thick] (q3) -- (out);
  \draw[thick] (clk.north) -| (q0.south);
  \draw[thick] (clk.north) -| (q1.south);
  \draw[thick] (clk.north) -| (q2.south);
  \draw[thick] (clk.north) -| (q3.south);
\end{tikzpicture}
\end{center}
\end{frame}
